\section{Multiplier, cost analysis, and proof of the main theorem}\label{sec:main}

This section assembles the complete integer multiplication algorithm from the fractional-shift bridge of
Section~\ref{sec:bridge} and the direct tensor FFT of Section~\ref{sec:direct-fft}, balances its parameters, and proves
Theorem~\ref{thm:main}. Two $n$-bit integers are cut into $(\log n)$-bit digits, packed into a single complex vector of
power-of-two length~$R$, and multiplied via two bridge calls, one pointwise squaring, nearest-integer rounding, and
carry propagation. Summing the normalized tape costs of all stages shows that each stage falls short of $n\log n$ by a positive power of
$\log n$. These five exponent gaps are collected in~\eqref{eq:gaps}, and the explicit parameter tuple of
Table~\ref{tab:walsh-rows} yields the saving exponent $\kappa=2^{-46}$ with strict positive slack. A first-time reader interested primarily in how the components fit together
and why $\kappa=2^{-46}$ emerges can read Sections~\ref{sec:main-alg}--\ref{sec:main-rows} and skim the symbolic cutoff
calculation in Section~\ref{sec:main-thm}.

\subsection{The algorithm and its root parameter family}\label{sec:main-alg}

This subsection defines the power-of-two grid dimensions as functions of the input bit length~$n$ and summarizes how the
top-level multiplier connects its subroutines.

\paragraph{Parameters of the root family.}
We consider two $n$-bit integers with $n$ at least the threshold $n_*$ defined in Section~\ref{sec:main-thm}; for
$n<n_*$, the machine invokes the independent inner multiplier directly. Let $b=\lfloor\log_2n\rfloor$ and
$N_{\rm dig}=\lceil n/b\rceil$, and let $R=2^k$ be the least power of two satisfying $R\ge2N_{\rm dig}$. Given fixed
rational exponents $\eps$ and $c$, chosen in Section~\ref{sec:main-rows}, define
\begin{equation}\label{eq:root-sizes}
  p=2b+2k+6,\qquad d=\lfloor k^\eps\rfloor,\qquad \ell=\lceil k/d\rceil,\qquad r=2^\ell,\qquad K=\lfloor d^c\rfloor.
\end{equation}
Because $0\le d\ell-k<d$, we set the first $d\ell-k$ axis lengths to $r/2$ and the remaining axis lengths, including
$t_d$, to $r$, so that $\prod_{i=1}^dt_i=R$. The alias width~$\tw$, internal working precision~$P$, grid precision
$q=P+2$, and component width $w=q+2$ follow the bridge schedule~\eqref{eq:schedule}; we call these sizes the \emph{root
family}. Fixed rational powers are computed by exact integer comparisons: for $\eps=\eps_1/\eps_2$, the integer
$d=\lfloor k^\eps\rfloor$ is the largest $y\in\N$ satisfying $y^{\eps_2}\le k^{\eps_1}$. Every transform length in this
root family is a power of two. As in Section~\ref{sec:direct-fft}, $\mathsf M(N)=N\,\mathsf m(N)$ denotes the cost of
the independent inner multiplier on $N$-bit integers.

\paragraph{Top-level pipeline.}
The multiplier packs the radix-$2^b$ digits of the two inputs into one complex vector of length~$R$ and recovers the
product polynomial from two one-dimensional transforms $F_R$ via the root-packing proposition
(Proposition~\ref{prop:root-packing}). Each $F_R$ is evaluated by the fractional-shift bridge under the bridge precision contract
(Theorem~\ref{thm:bridge-precision}); the axis exposures of Lemma~\ref{lem:bridge-D} bring each axis $a=d,\dots,2$
innermost for spreading with the negative half-bin carrier. Inside the bridge, the direct tensor FFT of
Corollary~\ref{cor:walsh-middle} evaluates the tensor transform. Its two address permutations, the layout map~$L$ and the
restoring map $U=(LB)^{-1}$ of Section~\ref{sec:direct-fft-layout}, are executed without error by the
positional-permutation theorem (Theorem~\ref{thm:positional-perm}). Finally, chronological deconvolutions
$B_2^{-1},\dots,B_d^{-1}$ undo the spreading to working precision.

The butterfly layer, the Walsh permutation routines, and the direct tensor FFT of
Sections~\ref{sec:walsh}--\ref{sec:direct-fft} assume that the dimension~$d$, the last-axis length~$r$, the chunk
width~$K$, and the grid precision stay inside fixed comparison bands. When $0<\eps<\frac1{14}$, $c>0$, and
$\eps(1+c)<1$, the next lemma verifies all of these size and growth conditions simultaneously for all sufficiently
large~$n$, both at the bridge grid precision~$q$ and at the lifted precision $q+8d$ of Lemma~\ref{lem:walsh-lift}.

\begin{lemma}[Uniform root family]\label{lem:root-family}
Fix any rational exponents $0<\eps<\frac1{14}$ and $c>0$ with $\eps(1+c)<1$ (in particular the primary parameter tuple
of Table~\ref{tab:walsh-rows}, as well as the comparison rows of Appendix~\ref{app:bluestein}). For all sufficiently
large~$n$: $RP=\Theta(n)$, $P=\Theta(q)=\Theta(k)=\Theta(\log n)$, $d=\Theta(P^\eps)$, $K=\Theta(P^{\eps c})$, and
$\ell=\Theta(P^{1-\eps})$. At both precisions $q'\in\{q,\ q+8d\}$, the comparison bands of
Proposition~\ref{prop:prop14-param} hold with the fixed constants
\begin{equation}\label{eq:bands}
  a_d=\tfrac14,\quad b_d=1,\quad a_r=\tfrac1{128},\quad b_r=\tfrac18,\quad C_\ell=1,\quad C_M=2,\quad C_w=2,\quad
  C_{\rm desc}=128,
\end{equation}
that is, $a_d(q')^{\eps}\le d\le b_d(q')^{\eps}$ and $2^{a_rq'/d}\le r\le2^{b_rq'/d}$. Moreover $dK=o(P)$, $K=o(\ell)$,
$K/\log P\to\infty$, and $r/P^j\to\infty$ for every fixed~$j$.
\end{lemma}

\emph{Proof idea.} From $R=\Theta(n/\log n)$ and $b=\Theta(\log n)$, the linear bridge schedule~\eqref{eq:schedule} pins
both $q$ and $q+8d$ within constant multiples of $k=\lg R$. Substituting $16k\le q'\le 64k$ and $\eps<1/14$ into the
definitions of $d$, $\ell$, and $K$ directly verifies each band constant in~\eqref{eq:bands} and the asymptotic growth
ratios.

\begin{proof}
\emph{Sizes.} From $2N_{\rm dig}\le R<4N_{\rm dig}$ we have $k\le b\le2k$ for large $n$ and $b/k\to1$, so the bridge
schedule~\eqref{eq:schedule} gives $P/k\to20$ and $q/k\to20$. In particular, $16k\le q\le32k$ for all $k\ge4096$, and
$q+8d\le64k$. Since $R=\Theta(n/\log n)$, we also have $RP=\Theta(n)$.

\emph{Bands.} Because $0<\eps<\frac1{14}$, once $k^\eps\ge2$ we have $\frac12k^\eps\le d\le k^\eps$. Since
$64^{1/14}=2^{3/7}<2$, for $q'\in\{q,q+8d\}$ we obtain
\[
  d\ge\tfrac12\Bigl(\frac{q'}{64}\Bigr)^\eps\ge\tfrac12\,64^{-1/14}(q')^\eps>\tfrac14(q')^\eps,
\]
while $d\le k^\eps\le(q')^\eps$. Because $k/d\le\ell\le2k/d$ and $16k\le q'\le64k$, we obtain
$q'/(128d)\le\ell\le q'/(8d)$. The remaining constants of~\eqref{eq:bands} hold as well: $d\ell\le k+d\le2k<q'$ and
$\lg(2\cdot R/r)\le k+1<2q'$. The nine-field headers have length $O(k)\le C_{\rm desc}q'$, and the component widths
$w=q+2$ and $w_*=q+2+8d$ satisfy $q'+2\le w_{q'}\le2q'$.

\emph{Growth.} Finally, $dK/k=O(k^{\eps(1+c)-1})\to0$, which implies $K=o(\ell)$ since $\ell=\Theta(k/d)$. The limits
$K/\log P\to\infty$ and $r/P^j\to\infty$ follow from $\eps c>0$ and $\ell\ge k^{1-\eps}$.
\end{proof}

\subsection{Exact recovery}\label{sec:main-recovery}

This subsection proves that two approximate bridge transforms and one rounded pointwise squaring determine every integer
convolution coefficient without error.

Packing the radix-$2^b$ digits of the two inputs as $x_j=2^{-b}(x^{(0)}_j+ix^{(1)}_j)$ lets the two-call
identity~\eqref{eq:two-call} express their cyclic convolution through two forward transforms $F_R$. The next proposition
tracks the bridge error bound~\eqref{eq:bridge-total} through the squaring and second transform to show that, at target
precision $p=2b+2k+6$, every scaled output entry lies within $\frac12$ of the true integer convolution coefficient.
Nearest-integer rounding and carry propagation then recover the exact $2n$-bit product.

\begin{proposition}[Two-call root packing]\label{prop:root-packing}
Let $x^{(0)}_j,x^{(1)}_j\in[0,2^b)$ ($0\le j<N_{\rm dig}$) be the radix-$2^b$ digits of the two $n$-bit inputs (with
$x^{(0)}_j=x^{(1)}_j=0$ for $N_{\rm dig}\le j<R$), and set $x_j=2^{-b}(x^{(0)}_j+ix^{(1)}_j)$ for $0\le j<R$. Compute
$\hat y\approx F_Rx$ by the bridge, $\hat w=\overline{\Quant_{p+2}(\hat y\odot\hat y)}$, and $\hat z\approx F_R\hat w$
by the bridge. Then
\[
  c_j=\mathrm{round}\bigl(-2^{2b-1}R^2\,\mathrm{Im}\,\hat z_j\bigr)
\]
equals the exact convolution coefficient $\sum_{u+v=j}x^{(0)}_u x^{(1)}_v$, and radix-$2^b$ carry propagation returns
the exact $2n$-bit product. Outside the two bridge calls, the cost is $O(R\,\mathsf M(P)+n)=O(RP\,\mathsf m(P)+n)$.
\end{proposition}

\emph{Proof idea.} Because $R\ge 2N_{\rm dig}$, the cyclic convolution of length~$R$ coincides with the non-cyclic digit
convolution. Applying the bridge error bound $0.177021\cdot 2^{-p}X$ of~\eqref{eq:bridge-total} twice, with
$X=\max(1,\|\cdot\|_\infty)$ evaluated on each bridge input, and accounting for the intermediate squaring and rounding
bounds the unscaled error by $1.2391\cdot 2^{-p}$. Multiplying by $2^{2b-1}R^2=2^{p-7}$ keeps the final error below
$1.2391\cdot 2^{-7}<\frac12$.

\begin{proof}
\emph{The identity.} Applying~\eqref{eq:two-call} to $2^{-b}x^{(0)}$ and $2^{-b}x^{(1)}$ gives
$x^{(0)}\circledast x^{(1)}=-2^{2b-1}R^2\,\allowbreak\mathrm{Im}\,\allowbreak F_R(\overline{(F_Rx)\odot(F_Rx)})$, which
has no cyclic wrap because the product polynomial has degree at most $2N_{\rm dig}-2<R$.

\emph{The error.} Put $u_0=2^{-p}$. Each $x_j$ lies on the $(p+2)$-grid with $\|x\|_\infty<\sqrt2$, so
Theorem~\ref{thm:bridge-precision} gives $\|\hat y-y\|_\infty<0.177021\sqrt2\,u_0<0.25035\,u_0$, where $y=F_Rx$ and
$\|y\|_\infty\le\|x\|_\infty<\sqrt2$. Since $p\ge14$, we have
$\|\hat y\odot\hat y-y\odot y\|_\infty\le(\|\hat y\|_\infty+\|y\|_\infty)\|\hat y-y\|_\infty<0.7082\,u_0$, and rounding
to $\Cg_{p+2}$ adds at most $\frac{\sqrt2}{8}u_0<0.17678\,u_0$. Hence $\|\hat w-\overline{y\odot y}\|_\infty<0.8850\,u_0$
and $\|\hat w\|_\infty<2.0001$. The second bridge call adds less than $0.177021\cdot2.0001\,u_0<0.35406\,u_0$, and
$\|F_R\|_\infty\le1$, so $\|\hat z-F_R(\overline{y\odot y})\|_\infty<1.2391\,u_0$. Multiplying by
$2^{2b-1}R^2=2^{2b+2k-1}=2^{p-7}$ yields an error below $1.2391\cdot2^{-7}<\frac12$, so rounding to the nearest integer
recovers every $c_j$ without error. On $\Cg_{p+2}$, if $n_j\in\Z$ is the numerator of $\mathrm{Im}\,\hat z_j$, then
$-2^{2b-1}R^2\,\mathrm{Im}\,\hat z_j=-n_j/512$.

\emph{Carries and cost.} Each coefficient satisfies $0\le c_j\le N_{\rm dig}(2^b-1)^2$. Setting $s_0=0$,
$\gamma_j=(c_j+s_j)\bmod2^b$, and $s_{j+1}=\lfloor(c_j+s_j)/2^b\rfloor$ produces the radix-$2^b$ digits $\gamma_j$ of
the product, which is strictly below $2^{2n}$, emitting $2n$ bits. Pointwise squaring uses $R$ complex multiplications
on $O(P)$-bit integers, and all remaining steps are linear scans.
\end{proof}

By Proposition~\ref{prop:root-packing}, the total running time $\mathsf M_{\rm W}(n)$ of the multiplier satisfies
\begin{equation}\label{eq:root-cost}
  \mathsf M_{\rm W}(n)\le2\,\mathsf T_{\rm W}(R,P)+O(R\,\mathsf M(P)+n),
\end{equation}
where $\mathsf T_{\rm W}(R,P)$ is the cost of one bridge call on the Walsh route, with no recursive call to
$\mathsf M_{\rm W}$ itself.

\subsection{Cost analysis of a bridge call}\label{sec:main-tables}

This subsection tallies all setup work, inner integer multiplications, and tape-routing steps in one bridge call,
normalizing each contribution by the stored bit volume $V=RP$.

Put $V=RP$. Because $\eps C_1<1$, all internal and external component widths, including the layer guard
$\Delta=\lceil C_0d^{C_1}\rceil=o(P)$ and the $8d$-bit Walsh lift, remain $O(P)$, so $V$ is the stored bit volume up to
constant factors. One-time and per-line or per-record setup work that is not repeated per coefficient covers cycle
decompositions, masks, headers, named-slot lists, filter taps, solver parameters, phase tables, and recursion
descriptors. Each such object is built in time polynomial in~$P$ across at most $(R/r+r)P^{O(1)}$ objects. Hence, for
fixed constants $C_{\rm set},D_{\rm set}$, the setup cost satisfies
\begin{equation}\label{eq:setup}
  \mathcal S(R,P)\le C_{\rm set}\bigl((R/r)P^{D_{\rm set}}+rP^{D_{\rm set}}+P^{D_{\rm set}}\bigr)=o(RP),
\end{equation}
by Lemma~\ref{lem:root-family} and $R/r\ge r/2$.

\paragraph{Operand widths of all inner integer products.}
Per-coefficient multiplications are delegated to an independent integer multiplier of cost
$\mathsf M(N)=N\,\mathsf m(N)$ on $N$-bit operands, assuming only $\mathsf m(N)=(\lg N)^{O(1)}$, as satisfied by
Sch\"onhage--Strassen~\cite{SS71} with $\mathsf m(N)=O(\lg N\lg\lg N)$. Its operands and output occupy dedicated work
tapes, leaving the caller's other tapes untouched and cleaning its workspace on return. Fixing a regular majorant
$\overline{\mathsf M}(N)=A\,N(\lg(2N))^J$ of $\mathsf M$ gives $\overline{\mathsf M}(c'N)=O(\overline{\mathsf M}(N))$
for every constant $c'>0$. Every ordinary integer product on the Walsh route occurs at one of four sites and has at most
$O(P^2)=O((\log n)^2)$ bits:
\begin{enumerate}[label=(\arabic*),leftmargin=*]
  \item \emph{Root packing} in Proposition~\ref{prop:root-packing}: $R$ complex squarings on $\Cg_{p+2}$, using
    $O(P)$-bit integer products of total cost $O(R\,\mathsf M(P))=O(V\,\mathsf m(P))$.

  \item \emph{Direct tensor FFT} in Theorem~\ref{thm:direct-fft} at $q=P+2$: $\ell$ last-axis levels and at most $\ell$
    record-axis combined twiddles, each performing four exact $(q+2)$-bit integer products per complex multiplication by
    Lemma~\ref{lem:disk-product}, of total cost $O(V\ell\,\mathsf m(P))$. Neither the butterfly layers of
    Proposition~\ref{prop:prop14-param} nor the positional permutations of Theorem~\ref{thm:positional-perm} perform
    general multiplications.

  \item \emph{Cyclic deconvolution} in part~(iv) of Lemma~\ref{lem:bridge-C}: $O(1)$ multiplications per entry on
    operands of at most $P+k+8=O(P)$ bits across $d-1$ axes, of total cost $O(Vd\,\mathsf m(P))$. Truncated binary
    powering on $O(P+k)$-bit integers for the closures in part~(v) of Lemma~\ref{lem:bridge-C} and rational-series
    evaluations on $O(P)$-bit integers in part~(iii) of Lemma~\ref{lem:bridge-C} and Lemma~\ref{lem:phase-table} belong
    to the setup term~\eqref{eq:setup}.

  \item \emph{Line spreading} in part~(ii) of Lemma~\ref{lem:bridge-C}: on each of $d-1$ axes, overlap-add blocks of
    $W_t=\min(2P+1,t_a)=\Theta(P)$ entries are packed via signed Kronecker substitution into $(4P+16)$-bit slots. Each
    block performs four exact products of integers of at most $W_t(4P+16)+1=O(P^2)$ bits across $O(R/P)$ blocks per
    axis, costing $O(Vd\,\mathsf m(P^2))$.
\end{enumerate}
In particular, because no polynomial-ring multiplication in $\Rr_r$ is performed, no $\Theta(rP)$-bit packed product
arises, and both $\mathsf m(P)$ and $\mathsf m(P^2)$ are $(\lg P)^{O(1)}$.

Under Theorem~\ref{thm:positional-perm}, both $L$ and $U=(LB)^{-1}$ cost
$O(V[(m_0/d+K)d^{\lambda'}+K])=O(V\ell d^{\lambda'})$. Combining this bound with Theorem~\ref{thm:direct-fft},
Lemma~\ref{lem:bridge-D}, and Lemma~\ref{lem:bridge-C} yields the normalized cost breakdown in
Table~\ref{tab:cost-walsh}.

\begin{table}[ht]
\centering
\small
\caption{Normalized running time of each stage in one bridge call on the Walsh route. Costs are divided by the bit
volume $V=RP$, with $d=\Theta(P^\eps)$, $\ell=\Theta(P^{1-\eps})$, and $K=\Theta(P^{\eps c})$; a row with gap $g$ has
normalized cost $O(P^{1-g})$ up to polylogarithmic factors in~$P$ where indicated
(Section~\ref{sec:main-gaps}).}\label{tab:cost-walsh}
\begin{tabular}{@{}lll@{}}
\toprule
operation & normalized cost & gap below $P^1$\\
\midrule
individual prefix butterfly rounds & $dK$ & $1-\eps(1+c)$\\
butterfly layers and positional layouts $L,U$ & $\ell d^{\lambda'}$ & $\eps(1-\lambda')$\\
complex twiddles and last-axis FFT & $\ell\,\mathsf m(P)$ & $\eps$, up to logarithms\\
axis exposures and restorations & $d(1+\ell^\tau)$ & $(1-\eps)(1-\tau)$\\
line spreading and cyclic solves & $d(\mathsf m(P)+\mathsf m(P^2))$ & $1-\eps$, up to logarithms\\
\bottomrule
\end{tabular}
\end{table}

Summing the rows of Table~\ref{tab:cost-walsh} and absorbing~\eqref{eq:setup} gives
\begin{equation}\label{eq:cost-W}
  \mathsf T_{\mathrm W}(R,P)=O\bigl(RP\,[dK+\ell d^{\lambda'}+\ell\,\mathsf m(P)+d(1+\ell^\tau)
  +d(\mathsf m(P)+\mathsf m(P^2))]\bigr).
\end{equation}
Here the direct axis exposures of Lemma~\ref{lem:bridge-D} contribute $d(1+\ell^\tau)$ in place of the
$d^2(1+\ell^\tau)$ adjacent-axis schedule of~\cite[\S9]{OAI26}. Executing $L$ and $U=(LB)^{-1}$ via
Theorem~\ref{thm:positional-perm} eliminates the external chunk-exchange term $PK^{\tau-1}$, while the internal packed
basis changes of Proposition~\ref{prop:prop14-param} are already included in $d^{\lambda'}$.

\subsection{Gaps and slack}\label{sec:main-gaps}

This subsection expresses the condition for saving a power of $\log n$ in terms of five exponent gaps. It also shows why
replacing chunk exchanges by Walsh permutations makes the exponent depend linearly, rather than quadratically, on the
network deficit introduced in Section~\ref{sec:overview}.

By Table~\ref{tab:cost-walsh}, a row with gap $g$ contributes $O(V\,P^{1-g}(\lg P)^{O(1)})$ steps. An admissible
parameter tuple therefore yields the saving exponent $\kappa$ whenever $0<\kappa<g_{\min}$, where $g_{\min}$ is the
minimum of the five gaps
\begin{equation}\label{eq:gaps}
  g_{\rm pre}=1-\eps(1+c),\quad g_{\rm lay}=\eps(1-\lambda'),\quad g_{\rm last}=\eps,\quad
  g_{\rm axis}=(1-\eps)(1-\tau),\quad g_{\rm short}=1-\eps.
\end{equation}
Any strict positive slack $s=g_{\min}-\kappa>0$ absorbs fixed powers of $\lg P$: for every integer $j\ge0$ and $P\ge2$,
\[
  (\lg(2P))^j\le4^j(\ln P)^j\le j!\,(4/s)^j\,P^{s},
\]
by the $j$-th term of the exponential series for $e^{s\ln P}$.

\paragraph{Linear vs.\ quadratic bottleneck.}
On the Walsh route, the binding gap is the butterfly-layer gap $g_{\rm lay}=\eps(1-\lambda')$. Since the layer
hypotheses require $\lambda'>\sigma$ and $\eps<1/C_1$, this gap is at most $(1-\sigma)/C_1$, which is \emph{linear} in
the complex network's exponent deficit $1-\sigma$. By contrast, if the layout $L$ is executed by width-$K$ chunk
exchanges as on a Bluestein middle transform in~\eqref{eq:old-bottleneck} of Appendix~\ref{app:bluestein}, the
normalized cost term $PK^{\tau-1}$ introduces the additional gap $g_{\rm chunk}=\eps c(1-\tau)$. Because
Proposition~\ref{prop:prop14-param} requires $c<\beta(1-\tau)/\tau$ and $\eps<1/C_1$, this forces
\[
  g_{\rm chunk}=\eps c(1-\tau)<\frac{\beta(1-\tau)^2}{C_1\tau}<\frac{(1-\tau)^2}{C_1\tau},
\]
capping $\kappa$ quadratically in $1-\tau$.

\subsection{Parameter choice}\label{sec:main-rows}

This subsection fixes the explicit rational parameter tuple for Theorem~\ref{thm:main} and verifies all admissibility
inequalities and positive cost gaps in exact rational arithmetic.

Proposition~\ref{prop:deficits} in Appendix~\ref{app:parameter-transfer} fixes the network exponents at $h=100$ and
$m=10^6$. With $\xi=2^{-42}$, the exact network counts of Fact~\ref{fact:oai-networks} together with the rational bound
$\ln 10^6<14$ certify that $\tau=1-4\xi=1-2^{-40}$ and $\sigma=1-\xi=1-2^{-42}$ satisfy $s_b/W_b<m^\tau$ and
$s_c/W_c<m^\sigma$. Likewise, the
bound $B_{\rm net}<m^{11}$, which gives $\nu_{\mathrm{net}}=11$ in~\eqref{eq:walsh-guard-constants}, certifies that the
guard exponent $C_1=\nu_{\mathrm{net}}+3=14$ is admissible in Proposition~\ref{prop:prop14-param}.
Table~\ref{tab:walsh-rows} records the parameter tuple used to prove Theorem~\ref{thm:main}.

\begin{table}[ht]
\centering
\caption{Explicit rational parameters for the primary Walsh multiplier. With $\xi=2^{-42}$, this tuple satisfies all
butterfly-layer and root-family constraints and yields the saving exponent $\kappa=2^{-46}$ in
Theorem~\ref{thm:main}.}\label{tab:walsh-rows}
\begin{tabular}{@{}lcccccccc@{}}
\toprule
$\kappa$ & $C_1$ & $\eps$ & $\beta$ & $c$ & $\tau$ & $\sigma$ & $\lambda$ & $\lambda'$\\
\midrule
$2^{-46}$ & $14$ & $1/15$ & $1/64$ & $\xi/256$ & $1-4\xi$ & $1-\xi$ & $1-63\xi/64$ & $1-31\xi/32$\\
\bottomrule
\end{tabular}
\end{table}

We verify the admissibility inequalities of Proposition~\ref{prop:prop14-param} and the gaps of~\eqref{eq:gaps} for
Table~\ref{tab:walsh-rows} in exact rational arithmetic:
\begin{itemize}[leftmargin=*]
  \item \emph{Layer and guard hypotheses.} We have $\max\{\tau,\sigma\}=1-\xi<\lambda=1-\frac{63}{64}\xi<1$ with margin
    $\lambda-\sigma=\frac{\xi}{64}>0$. Since $c/\beta=\xi/4$, we have
    $\tau(1+c/\beta)=(1-4\xi)(1+\xi/4)=1-\frac{15}{4}\xi-\xi^2$, giving the strict margin
    $\lambda-\tau(1+c/\beta)=\frac{177}{64}\xi+\xi^2>0$. Moreover $\sigma+\beta(1-\sigma)=1-\frac{63}{64}\xi=\lambda$, so
    $\max\{\lambda,\sigma+\beta(1-\sigma)\}=\lambda<\lambda'=1-\frac{31}{32}\xi<1$ with margin
    $\lambda'-\lambda=\frac{\xi}{64}>0$. Finally, $C_1=14$ is admissible and $1-\eps C_1=\frac1{15}>0$.

  \item \emph{Gaps and slack.} Evaluating~\eqref{eq:gaps} gives
    \[
      g_{\rm lay}=\frac{31}{480}\xi,\qquad g_{\rm axis}=\frac{56}{15}\xi,\qquad g_{\rm last}=\frac1{15},\qquad
      g_{\rm short}=\frac{14}{15},\qquad g_{\rm pre}=\frac{14}{15}-\frac{\xi}{3840}>\frac{13}{15}.
    \]
    The unique minimum is $g_{\min}=g_{\rm lay}=\frac{31}{480}\xi$. Comparing against $\kappa=2^{-46}=\frac{\xi}{16}
    =\frac{30}{480}\xi$ gives strict positive slack $s=g_{\min}-\kappa=\frac{\xi}{480}>0$, with ratio
    $g_{\min}/\kappa=31/30$.
\end{itemize}

The next remark isolates the architectural gain of the Walsh middle transform by evaluating the saving exponent under
the untuned network constants printed in~\cite{OAI26}.

\begin{remark}[Exponent with the printed constants of~\cite{OAI26}, and comparison rows]\label{rem:printed-optimal}
To isolate the gain of the Walsh middle transform from parameter retuning, fix the structural constants $\beta=\frac12$,
$\tau=\sigma=1-2^{-50}=1-2u$ ($u=2^{-51}$) and $C_1=20$ printed in~\cite[Prop.~14]{OAI26}. With these constants, the
printed statements~\cite[Lemma~12, Prop.~14]{OAI26} apply directly without Appendix~\ref{app:parameter-transfer}. Taking
$\eps=1/21$, $c=u/4$, $\lambda=1-129u/128$ and $\lambda'=1-63u/64$ satisfies all hypotheses of~\cite[Prop.~14]{OAI26}
(with $\tau(1+c/\beta)=1-\frac32u-u^2<\lambda$, $\sigma+\beta(1-\sigma)=1-u<\lambda'$, and $1-\eps C_1=1/21$) and gives
$g_{\min}=g_{\rm lay}=3u/64=3\cdot2^{-57}$, yielding $\kappa=u/32=2^{-56}$ with slack $2^{-57}$. Because any choice of
the free parameters with these fixed structural constants must satisfy $\lambda'>\sigma+\beta(1-\sigma)=1-u$ and
$\eps<1/C_1=1/20$, one has $\kappa<g_{\rm lay}=\eps(1-\lambda')<u/20=\frac85\,2^{-56}<2^{-55}$; thus $2^{-56}$ is the
largest exponent of the form $2^{-j}$ ($j\in\N$) that our analysis admits with those constants.
Appendix~\ref{app:bluestein} gives the full comparison with a Bluestein middle transform (reproducing $2^{-182}$ at the
parameters of~\cite[(38)]{OAI26} and reaching $2^{-84}$ with retuned parameters) together with the $C_1=20$ Walsh row
($\kappa=2^{-47}$).
\end{remark}

\subsection{The cutoff and proof of Theorem~\ref{thm:main}}\label{sec:main-thm}

This subsection defines the symbolic input-length threshold $n_*$ above which all asymptotic bounds activate
simultaneously, and completes the proof of Theorem~\ref{thm:main}.

Because the tape procedures derived from Facts~\ref{fact:oai-streams}--\ref{fact:oai-networks} have unspecified implied
constants and eventual thresholds, we define the fallback threshold symbolically; Section~\ref{sec:discussion} discusses
the resulting astronomical activation sizes.

\paragraph{Cutoff parameters.}
For the tuple of Table~\ref{tab:walsh-rows}, put $s=g_{\min}-\kappa=\xi/480>0$. Fix integers $C\ge2$, $j\ge1$, and
$D_{\rm set}\ge1$ such that every normalized cost term in~\eqref{eq:cost-W} and~\eqref{eq:root-cost} is at most
$CP^{1-g_{\min}}(\lg(2P))^j$ and~\eqref{eq:setup} holds with coefficient at most~$C$ and degree at most $D_{\rm set}$.
Let $P_{\rm tape}$ be a common precision threshold under the bands~\eqref{eq:bands} for both the butterfly layer of
Proposition~\ref{prop:prop14-param} and the arbitrary-width interchange of Lemma~\ref{lem:transfer12}. Let $C_0$ be the
guard constant~\eqref{eq:walsh-guard-constants} of Proposition~\ref{prop:prop14-param} from
Appendix~\ref{app:parameter-transfer}. For $\alpha_0>0$ and $b_*\ge1$ put
$\Phi(\alpha_0,b_*)=2^{\lceil(64b_*/\alpha_0)^2\rceil}$, and put
$K_{\log}=2^{\lceil\max\{2,(16j/s)^2,(8/s)\lg(2C)\}\rceil}$. Define the constant
\begin{equation}\label{eq:cutoff}
\begin{aligned}
  K_*=\Bigl\lceil\max\Bigl\{&4096,\ P_{\rm tape},\ (4C_0)^{1/(1-\eps C_1)},\ \Phi(\eps c,1),\ 2^{2/\eps},\
  2^{2/(1-\eps(1+c))},\\
  &\Phi\bigl(1-\eps,\lceil D_{\rm set}+2+\lg(8C)\rceil\bigr),\ K_{\log}\Bigr\}\Bigr\rceil
\end{aligned}
\end{equation}
and set $n_*=2^{4K_*}$; here $K_*$ is a fixed threshold rather than the chunk width~$K$.

For $n\ge n_*$ we have $b=\lfloor\log_2n\rfloor\ge4K_*$, so $2N_{\rm dig}\ge2n/b>2^{2K_*}$ and hence $k\ge K_*$. We now
verify that $k\ge K_*$ activates all asymptotic bounds simultaneously. If $k\ge\Phi(\alpha_0,b_*)$ and $z=\log_2k$, then
$z\ge(64b_*/\alpha_0)^2$ implies $(\ln2)\alpha_0z\ge32b_*\sqrt z\ge\ln(b_*(z+5)^2)$, so $k^{\alpha_0}\ge b_*(z+5)^2$.
Since $q\le32k$ gives $\log_2q\le z+5$, taking $\alpha_0=\eps c$ and $b_*=1$ makes $K\ge6$ and $K/\log P$ large. Taking
$\alpha_0=1-\eps$ and $b_*=\lceil D_{\rm set}+2+\lg(8C)\rceil$ gives $r\ge8CP^{D_{\rm set}+1}$ and
$R/r\ge r/2\ge4CP^{D_{\rm set}+1}$, so the setup~\eqref{eq:setup} is at most $RP$. The term $(4C_0)^{1/(1-\eps C_1)}$
ensures $\Delta\le C_0k^{\eps C_1}+1\le k/2$, and the terms $2^{2/\eps}$ and $2^{2/(1-\eps(1+c))}$ ensure $d\ge2$ and
$dK\le k/4$. Finally, $K_{\log}$ ensures $\ln C+j\ln\lg(2P)\le\frac s2\ln P$, so half of the slack $s$ survives all
implied constants and polylogarithmic factors.

\emph{Proof idea.} Below the fixed threshold $n_*$, the machine delegates multiplication directly to a standard
$N(\log N)^{O(1)}$ inner multiplier in bounded time. For $n\ge n_*$, the cutoff~\eqref{eq:cutoff} makes every size band,
guard-bit margin, and setup bound derived in Sections~\ref{sec:main-alg}--\ref{sec:main-rows} hold simultaneously. Two
calls to the fractional-shift bridge, instantiated with the direct complex tensor FFT and exact Walsh-layer address
permutations, recover the exact $2n$-bit product by Proposition~\ref{prop:root-packing}. The strict slack
$s=g_{\min}-\kappa=\xi/480>0$ absorbs all polylogarithmic factors in $P=\Theta(\log n)$, yielding total step count
$O(n(\log n)^{1-2^{-46}})$ on a single uniform multitape machine.

\begin{proof}[Proof of Theorem~\ref{thm:main}]
Fix the parameter tuple of Table~\ref{tab:walsh-rows} with $\kappa=2^{-46}$ and the cutoff $n_*=2^{4K_*}$
from~\eqref{eq:cutoff}. Let the independent inner multiplier have cost $\mathsf M(N)=N(\log N)^{O(1)}$, such as
Sch\"onhage--Strassen~\cite{SS71}, and for $n<n_*$ invoke this inner multiplier directly on the two $n$-bit inputs. We
verify parameter admissibility, correctness, outer CRT-freeness, running time, and machine uniformity for $n\ge n_*$.

\emph{Admissibility.} Proposition~\ref{prop:deficits} certifies the network exponents $\tau=1-2^{-40}$ and
$\sigma=1-2^{-42}$ and the admissibility of $C_1=14$ at $h=100$, and the strict margins of
Proposition~\ref{prop:prop14-param} hold as checked after Table~\ref{tab:walsh-rows}. Since $\eps=1/15<1/14$ and
$\eps(1+c)<1$, Lemma~\ref{lem:root-family} guarantees the comparison bands~\eqref{eq:bands}, component widths, complete
rectangular shapes, and long-record conditions at both precisions $q=P+2$ and $q_*=q+8d$.

\emph{Correctness and outer structure.} All transform lengths $R=2^k$ and $t_1,\dots,t_d\in\{r/2,r\}$ are powers of two;
the outer reduction uses no Chinese remainder map, no prime transform length, no coprime-length resampling, and no
Fourier-grid expansion. Inside each bridge call, Theorem~\ref{thm:positional-perm} executes the layout map $L$ and the
restoring map $U=(LB)^{-1}$ without error on $\Cg_q$, returning from the lifted precision $q+8d$ to $\Cg_q$ at proved
exact permutation boundaries. Theorem~\ref{thm:direct-fft} and Corollary~\ref{cor:walsh-middle} therefore supply a
middle transform at grid precision~$q$ with $E_{\rm mid}=12\ell$ and $\delta_{\rm mid}=0$. By
Theorem~\ref{thm:bridge-precision} and Lemma~\ref{lem:bridge-D}, the two bridge transforms satisfy the precision
bound~\eqref{eq:bridge-total}, and Proposition~\ref{prop:root-packing} recovers the exact $2n$-bit product.

\emph{Inner products and running time.} As verified in Section~\ref{sec:main-tables}, every ordinary integer product in
the algorithm has at most $O(P^2)=O((\log n)^2)$ bits and is executed by the independent multiplier $\mathsf M$, with no
recursive call to the algorithm of Theorem~\ref{thm:main}. By~\eqref{eq:root-cost}, \eqref{eq:cost-W}, the logarithmic
absorption of Section~\ref{sec:main-gaps}, and the positive slack $s=g_{\min}-\kappa=\xi/480>0$, the running time for
$n\ge n_*$ satisfies $\mathsf M_{\rm W}(n)=O(RP\cdot P^{1-\kappa})$. Since $RP=\Theta(n)$ and $P=\Theta(\log n)$ by
Lemma~\ref{lem:root-family}, this is $O(n(\log n)^{1-2^{-46}})$ uniformly for $n\ge n_*$. On the finitely many input
lengths $n<n_*$, the direct call to the inner multiplier takes bounded time, absorbed into the implied constant.

\emph{Uniformity.} All lists of axis lengths, masks, phase and filter tables, and positional-map decompositions are
generated deterministically by the machine as tape data. Moreover, every subroutine is a sequential composition of
fixed-tape procedures over a fixed finite alphabet and tape count independent of~$n$.
\end{proof}

Finally, the next remark explains how the black-box inner multiplier $\mathsf M$ is instantiated when one also wishes to
avoid prime transform lengths inside the $O((\log n)^2)$-bit scalar products, as in our Lean~4 formalization.

\begin{remark}[CRT-free inner multiplication]\label{rem:crt-free}\label{cor:crt-free}
Theorem~\ref{thm:main} treats the inner multiplier $\mathsf M$ as a black box on $O((\log n)^2)$-bit integers and does
not depend on~\cite{MS}. If one instantiates $\mathsf M$ instead with the two-dimensional fractional-shift $O(N\log N)$
multiplier of~\cite{MS} in place of~\cite{SS71}, then the inner integer products likewise use only power-of-two Fourier
transforms without a Chinese remainder map or prime transform lengths. Its written
analysis uses faithful phase and filter tables as in Section~\ref{sec:bridge}. In the formal proof of Theorem~\ref{thm:lean},
the inner multiplier is supplied by an independent Lean~4 formalization of the two-dimensional multiplier of~\cite{MS}
(headline \texttt{IntMulTM.\allowbreak intMul\_\allowbreak TM2\_\allowbreak nlogn}), developed for this paper and
kernel-checked as part of~\cite{Lean26} (Section~\ref{sec:lean}). The companion manuscript~\cite{MS} distributes its
own, separate Lean development (headline \texttt{IntMul.\allowbreak tmIntMul\_\allowbreak evalsToInTime\_\allowbreak n\_\allowbreak log\_\allowbreak n}).
\end{remark}
